\chapter*{Epilogue: One Diagram of the Field}

\addcontentsline{toc}{chapter}{Epilogue: One Diagram of the Field}

Looking back at the 38 entries, the intellectual architecture can be
drawn as one diagram:

\begin{center}
\begin{tabular}{ccc}
Wilson 1974 & $\longrightarrow$ & HMC / smearing / distillation (Ch.~2)\\[2pt]
$\downarrow$ & & $\downarrow$\\[2pt]
gradient flow program (Ch.~3) & & LaMET / pseudo-PDF revolution (Ch.~4)\\[2pt]
$\downarrow$ & & $\downarrow$\\
\multicolumn{3}{c}{renormalization battle: RI/MOM, auxiliary fields,
hybrid schemes (Ch.~5)}\\[-1pt]
\multicolumn{3}{c}{$\downarrow$}\\[-1pt]
\multicolumn{3}{c}{gluon PDFs: flowed and boosted lines converge (Ch.~6)}\\[2pt]
$\swarrow$ & & $\searrow$\\[2pt]
global fits: CT18, NNPDF (Ch.~7) & & stochastic quantization reborn:
AI sampling (Ch.~8)
\end{tabular}
\end{center}

Three observations close this selection.

\textbf{First}, the field's decisive theoretical events were cheap in
machinery but expensive in thought. The quasi-PDF proposal is a one-page
definition; the factorization theorem is an OPE argument; the flow's
finiteness is a power-counting proof. What made them decisive was that
each converted a conceptual impossibility --- ``the lattice cannot reach
light cone $x$-dependence'', ``composite operators need infinite
counterterms'' --- into a finite, computable program.

\textbf{Second}, convergence beats competition. Quasi- and pseudo-PDFs,
flowed and boosted operators, Hessian and Monte Carlo fits: every pair of
rivals turned out to be two coordinates on one object, and the
cross-checks between them are now the field's main systematics tool.

\textbf{Third}, this corpus is not closed. Its newest layer ---
agentic workflows that orchestrate the deterministic pipeline documented
in these chapters --- treats the papers above as domain logic to be
automated rather than replaced. Whether that layer earns a chapter of its
own in a future edition depends on precisely the standards this selection
was compiled under: reproducibility, error control, and physics first.
